\documentclass[11pt,a4paper]{Plantilla}
\usepackage[utf8]{inputenc}
\usepackage[spanish]{babel}
\usepackage{newtxtext,newtxmath}
\usepackage{setspace}
\singlespacing%
\usepackage{listings}
\usepackage{xcolor}
\usepackage{natbib}
\usepackage{bibentry}
\usepackage[table]{xcolor}
\usepackage{tabularx}
\usepackage{float}
\usepackage{array}
\usepackage{colortbl}
\usepackage{caption}
\usepackage{booktabs}
\usepackage{verbatim}
\usepackage{multicol}
\setlength{\parindent}{0pt}
\setlength{\parskip}{5pt}
\usepackage{url}
\nobibliography*
\definecolor{cabecera}{RGB}{208,226,243}
\asignatura{Complementos de Ingeniería de Software}
\titulo{Hito 1}
\prof{Daniel Moreno}
\alumnos{Christopher Abarca\\
Benjammin Echeverria\\
Maico Huillca \\
Camila Romero}
\begin{document}
\fairemarges%
\portada%

\section*{Información preliminar}


\subsection*{Equipo de Desarrollo}
\begin{itemize}
    \item Christopher Abarca (cabarcagutierrez@uft.edu)
    \item Benjammin Echeverria (becheverriav@uft.edu)
    \item Maico Huillca (mhuillcaq@uft.edu)
    \item Camila Romero (cromeroc2@uft.edu)
\end{itemize}

\subsection*{Contraparte}
\begin{itemize}
    \item Daniel Moreno (dmoreno@uft.cl)
\end{itemize}


\subsection*{Historia de versiones}

\begin{table}[H]
\centering
\caption{Historial de versiones}
\label{tab:historial-versiones}
\begin{tabular}{|c|c|c|c|}
\hline
\textbf{Versión} & \textbf{Fecha} & \textbf{Descripción} & \textbf{Autor} \\
\hline
1.0 & 2026-04-05 & Versión inicial & Christopher Abarca \\
\hline
1.1 & 2026-04-11 & Entrega 1 & Todos\\
\hline
\end{tabular}
\end{table}
\newpage

\tableofcontents
\listoftables
\newpage

\section{Introducción}%(maico)

El proyecto Finis Trabaja surge como una iniciativa estratégica de la Dirección de Vinculación con el Medio de la Universidad Finis Terrae para centralizar y modernizar el vínculo profesional entre su comunidad estudiantil y el mercado laboral. Históricamente, esta relación ha dependido de mecanismos informales e ineficientes, tales como correos masivos, carteles en los pasillos y grupos de mensajería que no garantizan la llegada efectiva de la información relevante a los interesados. Ante esta problemática, se propone el desarrollo de una plataforma web integral que funcione como un ecosistema digital eficiente para el intercambio de oportunidades, permitiendo que empresas externas y áreas internas de la institución gestionen vacantes de manera estructurada mientras los estudiantes acceden a ofertas alineadas con sus perfiles académicos. La relevancia de esta intervención radica en la creación del Currículum Vitae Virtual (CVV), una herramienta que permite a los alumnos gestionar su identidad profesional, especialidades y preferencias laborales de forma dinámica. El componente más innovador del sistema es su módulo de recomendación automática, el cual cruza datos de las vacantes con los perfiles CVV para identificar grados de compatibilidad basados en la carrera, pretensiones salariales, jornada y modalidad de trabajo, optimizando así la transición de los estudiantes hacia la práctica profesional.


\subsection{Contexto operacional}%(benja)

La Universidad Finis Terrae se constituye como una institución de educación superior compleja, integrada por diversas facultades y carreras técnicas y profesionales. En este ecosistema académico, la conexión y los flujos de comunicación institucional deben actuar como un puente estratégico fundamental entre el estudiantado y su inserción en el mercado laboral. No obstante, actualmente se identifica que esta comunicación no opera en su totalidad de forma eficiente, encontrándose dispersa y fragmentada, lo que impide que la información crítica alcance a los perfiles que la requieren y compromete el cumplimiento de los objetivos de empleabilidad de la institución.

\subsubsection{Escenario actual y problemática}

Históricamente, la relación entre la oferta de prácticas profesionales o empleos y la demanda estudiantil se ha gestionado a través de medios informales y descentralizados. Estos mecanismos no poseen capacidades de filtrado ni garantizan que la información sea recibida por los perfiles académicos adecuados. Entre las prácticas comunes se encuentra el uso de correos electrónicos masivos, los cuales suelen ser clasificados como spam por los servidores, la fijación de carteles físicos en las dependencias de la universidad y la difusión de vacantes mediante grupos de mensajería instantánea como WhatsApp o Telegram.

Esta desorganización informativa genera una saturación de datos irrelevantes para el usuario: es común que a un estudiante le llegue información de áreas ajenas a su competencia, mientras que las ofertas pertinentes no alcanzan su perfil por falta de exposición dirigida. Esta situación no solo limita las opciones de desarrollo del estudiante, sino que genera una percepción de falta de profesionalismo ante las empresas externas interesadas en el talento de la universidad, debilitando el vínculo entre la academia y el sector productivo.

\subsubsection{Actores en el entorno operativo}

El contexto operativo de la plataforma denominada ``Finis Trabaja'' se centra en la interacción sinérgica de tres grupos de interés, cada uno con necesidades operacionales específicas:

\begin{itemize}
    \item \textbf{Estudiantes y Egresados:} Este grupo constituye la base de usuarios postulantes que buscan visibilizar su perfil académico en un mercado laboral altamente competitivo. Ante la actual necesidad de postular de manera manual y reiterativa, surge el concepto de \textit{Currículum Vitae Virtual} (CVV). Este perfil digital no solo centraliza competencias y certificados, sino que permite declarar variables críticas como pretensiones económicas, disponibilidad horaria y preferencia de modalidad (presencial, híbrida o remota), permitiendo descartar automáticamente ofertas que no se ajusten a sus aspiraciones.
    
    \item \textbf{Empresas y Áreas Internas de la Universidad:} Representan el sector demandante que requiere un canal formal y estructurado para la publicación y gestión de vacantes. Actualmente, estas entidades dependen de procesos manuales y tediosos para la difusión de ofertas a través de múltiples canales no integrados. El entorno operativo requiere que estas empresas puedan gestionar el ciclo de vida de sus ofertas, detallando con precisión los roles, sueldos, tipos de jornada y modalidades de trabajo.
    
    \item \textbf{Administradores:} Actores clave encargados de la gobernanza del sistema. Su rol operativo se centra en la supervisión, validación de empresas que solicitan ingreso al portal y la auditoría de las ofertas publicadas para asegurar que cumplan con los estándares institucionales.
\end{itemize}

\subsubsection{Infraestructura y disponibilidad}

La plataforma Finis Trabaja debe operar como un sistema web centralizado, accesible desde internet y utilizable desde computadores y dispositivos móviles. Debe ser compatible con navegadores modernos y permitir la consulta de ofertas, actualización de perfiles y postulaciones en distintos horarios. Aunque en esta etapa no se define integración obligatoria con sistemas institucionales específicos, la solución debe contemplar disponibilidad continua, acceso remoto y mecanismos básicos de respaldo de información.
\subsection{Propósito del sistema}%(benja)

El proyecto "Finis Trabaja" se crea principalmente para solucionar el desorden que hay hoy en día y lograr una conexión directa entre los alumnos y las oportunidades de trabajo, ya sean de la misma universidad o de empresas de afuera. La idea es que el sistema no sea solo un lugar para colgar anuncios, sino una herramienta que ayude de verdad a que el estudiante encuentre algo que le sirva.

Para que el sistema funcione y cumpla su objetivo, se enfoca en estos puntos:

\begin{enumerate}
    \item \textbf{Tener todo en un solo lugar:} El propósito es que ya no dependamos de WhatsApp o de correos que nadie lee. Al tener una plataforma oficial, las empresas pueden publicar sus vacantes con todos los detalles (sueldo, horario, si es remoto o no) y los alumnos saben exactamente dónde ir a buscar.
    
    \item \textbf{Mejorar el perfil del alumno (CVV):} Queremos que el estudiante pueda mostrar más que un simple papel. Con el Currículum Vitae Virtual, el alumno puede poner qué es lo que realmente le interesa, cuánto quiere ganar y en qué horarios puede trabajar. Así, el sistema ayuda a que no pierda tiempo viendo ofertas que no tienen nada que ver con lo que busca.
    
    \item \textbf{Conectar de forma automática:} Lo más importante del software es el módulo de recomendación. El propósito es que el sistema haga el trabajo pesado de revisar quién califica para qué puesto. Así, la universidad puede ver rápido qué alumnos son los mejores para cada oferta según su carrera y lo que ellos mismos pusieron en su perfil.
    
    \item \textbf{Dar más visibilidad:} Al final, lo que se busca es que el alumno esté más cerca de conseguir un trabajo que le guste y que las empresas vean el talento que hay en la Finis Terrae de una forma más profesional y organizada.
\end{enumerate}

Este software se hace para que conseguir trabajo o práctica sea mucho más fácil y rápido para todos, pasando de un sistema de avisos sueltos a una plataforma que conecta a la gente por sus capacidades e intereses reales.




\subsubsection{Objetivo general}%(cami)
\textbf{Objetivo general del proyecto:} Desarrollar una plataforma web que centralice la gestión de oportunidades laborales y facilite la vinculación entre estudiantes, egresados, universidad y entidades oferentes.

\textbf{Objetivo general de la entrega:} Analizar el contexto del problema y documentar de forma inicial la propuesta de solución, definiendo propósito, alcance y objetivos del proyecto.
\subsubsection{Objetivos específicos}%(cami)
\textbf{Objetivos específicos del proyecto:}
\begin{itemize}
    \item Identificar las necesidades de estudiantes, empresas y administradores.
    \item Diseñar una solución que permita publicar ofertas, gestionar perfiles y facilitar postulaciones.
    \item Incorporar un mecanismo de recomendación entre vacantes y perfiles.
\end{itemize}

\textbf{Objetivos específicos de la entrega:}
\begin{itemize}
    \item Describir el contexto operacional y la problemática actual.
    \item Definir el propósito del sistema y el alcance preliminar del proyecto.
    \item Establecer una base documental inicial.
\end{itemize}
\subsection{Alcance}%(cami)
La solución del proyecto conlleva cumplir con las necesidades de los implicados y con ello amainar su incertidumbre. En este proyecto, los usuarios corresponden a estudiantes, empresas y administradores. 

\vspace{5pt}

En un aspecto generalizado, el proyecto está dirigido a los estudiantes de la Universidad Finis Terrae que buscan ingresar al área laboral, donde puedan aplicar conocimientos de sus respectivas carreras. Más específicamente, busca conectar a los estudiantes con trabajos que cumplan sus especificaciones y viceversa, de modo que, por medio de filtros, se disminuya el tiempo de búsqueda.

\vspace{5pt}

Respecto a las empresas, estas cuentan con distintos medios y plataformas para darse a conocer. Sin embargo, si no llegan a las personas correctas, podrían perder la oportunidad de localizar empleados que cumplan con sus requerimientos. La solución propone mejorar la comunicación con estudiantes y egresados, facilitando la identificación de perfiles que cumplan con las competencias necesarias.

\vspace{5pt}

Los administradores forman parte del sistema y actúan como intermediarios entre estudiantes y empresas, asegurando que la información sea confiable y correcta. Por lo tanto, la solución proporciona un canal que recibe y redistribuye la información de forma ordenada.

\vspace{5pt}

Finalmente, el sistema contempla funcionalidades como la publicación de ofertas laborales, el filtrado de información, la gestión de perfiles y un mecanismo de recomendación que permite relacionar perfiles de estudiantes con ofertas laborales, facilitando una conexión más eficiente entre ambas partes. No obstante, no considera procesos externos como la contratación directa ni la gestión fuera del entorno de la plataforma.

\subsection{Definiciones, acrónimos y abreviaturas}%(cami)
Estudiantado: Conjunto de estudiantes que forman parte de la comunidad académica de la Universidad Finis Terrae, quienes buscan oportunidades laborales relacionadas con sus áreas de estudio.

\vspace{5pt}

Egresados: Personas que han completado sus estudios en la Universidad Finis Terrae y están en búsqueda de oportunidades laborales para aplicar sus conocimientos y habilidades adquiridos durante su formación académica.

\vspace{5pt}

Mercado laboral: Conjunto de oportunidades de empleo disponibles para los estudiantes y egresados de la Universidad Finis Terrae.

\vspace{5pt}

Desentralizado: Sistema o proceso que no está centralizado, es decir, que no tiene un punto único de control o gestión.

\vspace{5pt}

Perfil academico: Conjunto de información que describe la formación académica, habilidades y experiencia de un estudiante o egresado.


\vspace{5pt}

Oferta laboral: Anuncio o publicación que describe una oportunidad de empleo, incluyendo requisitos, responsabilidades y beneficios.

\vspace{5pt}

Usuario: Persona que interactúa con la plataforma, ya sea como estudiante, empresa o administrador.

\vspace{5pt}

CVV: Currículum Vitae, documento que resume la formación académica, experiencia laboral y habilidades de una persona.

\vspace{5pt}

Entorno de desarrollo: Conjunto de herramientas, tecnologías y metodologías utilizadas para el desarrollo del sistema.

\vspace{5pt}

Gobernanza del sistema: Conjunto de políticas, procedimientos y estructuras organizativas que aseguran la gestión efectiva y el cumplimiento de los objetivos del sistema.

\vspace{5pt}

Administradores: Personas encargadas de gestionar y supervisar el funcionamiento de la plataforma, asegurando que la información sea confiable y correcta.

\vspace{5pt}

Multiplataforma: Capacidad de la plataforma para ser accesible y funcional en diferentes dispositivos y sistemas operativos.

\vspace{5pt}

Alta disponibilidad: Característica del sistema que garantiza su funcionamiento continuo y sin interrupciones, minimizando el tiempo de inactividad.

\vspace{5pt}

Arquitectura web: Estructura y diseño del sistema que permite su funcionamiento a través de internet, facilitando el acceso y la interacción de los usuarios con la plataforma.

\vspace{5pt}

Tiempo de respuesta: Tiempo que tarda el sistema en procesar una solicitud y proporcionar una respuesta al usuario.

\vspace{5pt}

Sistema: Conjunto de elementos interrelacionados que trabajan juntos para lograr un objetivo común, en este caso, la plataforma web que conecta a estudiantes y egresados con oportunidades laborales.

\vspace{5pt}

Software: Conjunto de programas, datos y documentación que forman parte del sistema, permitiendo su funcionamiento y uso por parte de los usuarios.

\vspace{5pt}

Modulos: Componentes o partes del sistema que cumplen funciones específicas, como la gestión de perfiles, publicación de ofertas laborales, filtrado de información y recomendación de oportunidades.

\vspace{5pt}

Plataforma web: Aplicación accesible a través de internet que permite a los usuarios interactuar con el sistema y acceder a sus funcionalidades.

\vspace{5pt}

Requisitos: Condiciones o capacidades que el sistema debe cumplir para satisfacer las necesidades de los usuarios y lograr los objetivos del proyecto.


\subsection{Referencias}%(chris)

\begin{thebibliography}{9}

\bibitem{linkedin}
LinkedIn Corporation,
\textit{LinkedIn --- Red profesional y bolsa de empleo},
2026.
[En línea]. Disponible en: \url{https://www.linkedin.com}.
[Accedido: 12 abr. 2026].

\bibitem{computrabajo}
Computrabajo S.A.,
\textit{Computrabajo --- Bolsa de trabajo y empleo en Chile},
2026.
[En línea]. Disponible en: \url{https://www.computrabajo.cl}.
[Accedido: 12 abr. 2026].


\bibitem{indeed}
Indeed, Inc.,
\textit{Indeed --- Bolsa de trabajo: búsqueda de empleo},
2026.
[En línea]. Disponible en: \url{https://cl.indeed.com}.
[Accedido: 12 abr. 2026].

\bibitem{bne}
Gobierno de Chile, Ministerio del Trabajo y Previsión Social,
\textit{Bolsa Nacional de Empleo (BNE)},
2026.
[En línea]. Disponible en: \url{https://www.bne.cl}.
[Accedido: 12 abr. 2026].



\end{thebibliography}


\section{Descripción de la solución}%(benja)


\subsection{Características de los participantes}%(benja)

\subsubsection{Stakeholders}%(benja)

\subsubsection{Usuarios del sistema}%(benja)

\subsection{Perspectiva del Producto según los Usuarios y Clientes}%(benja)


\subsection{Ambiente Operacional de la Solución}%(cami)

\subsubsection*{Estándares}
El sistema a desarrollar se basa en una arquitectura web, que funcionará con módulos específicos para la gestión de perfiles de estudiantes (CVV), publicación de ofertas laborales por parte de empresas, filtrado de información y un sistema de recomendación que cruce datos entre vacantes y perfiles.

Por lo tanto el sistema debe contar con mecanismos de seguridad para proteger la información personal y garantizar la privacidad de los usuarios. Las leyes de protección de datos vigentes 21719 y 21720, implementan medidas de seguridad como cifrado de datos, autenticación robusta y control de acceso para evitar el uso indebido de la información. (Del Congreso Nacional, s. f.) Lo que ayuda a cumplir con los requisitos de privacidad y seguridad establecidos por la legislación chilena, protegiendo los datos personales de los usuarios.

Para aplicar las leyes 17719 y 21720, el sistema debe implementar mecanismos para que los usuarios puedan ejercer sus derechos de acceso, rectificación, cancelación y oposición (ARCO) sobre sus datos personales. Esto implica proporcionar opciones para que los usuarios puedan actualizar su información, eliminar su cuenta o solicitar la restricción del procesamiento de sus datos en caso de ser necesario. 

Por otro lado existen las normas de seguridad de la información en informática, ISO/IEC 27001, que establece un marco de referencia para la gestión de la seguridad de la información. Lo que incluye la implementación de controles de seguridad, la evaluación de riesgos y la mejora continua del sistema de gestión de seguridad de la información (SGSI). (International Organization for Standardization, s. f.). Adicionalmente esta la norma ISO/IEC 25010, que proporciona un marco para la evaluación de la calidad del software, incluyendo aspectos como la funcionalidad, la fiabilidad, la usabilidad, la eficiencia, la mantenibilidad y la portabilidad (International Organization for Standardization, s. f.).

Estas normas se pueden aplicar en el desarrollo del sistema a través de controles de seguridad, pruebas de calidad y evaluación del software para garantizar que cumpla con los estándares de seguridad y calidad establecidos por la industria.

Considerando los aspectos de seguridad, calidad y protección de datos mencionados anteriormente, es necesario definir una base tecnológica que permita implementar estos requerimientos dentro de la solución propuesta.

\subsubsection*{Tecnología}
En el ámbito tecnológico, el proyecto se desarrollará en base al lenguaje de programación de python. Por lo tanto el desarrollo web se realizará utilizando frameworks de python como Django o Flask, que permiten la creación de aplicaciones web robustas y escalables. En este caso Django, un framework de alto nivel que sigue el patrón de diseño Modelo-Vista-Controlador (MVC) y proporciona una estructura sólida para el desarrollo de aplicaciones web, es elegido para el proyecto debido a la magnitud del mismo y las necesidades estipuladas en los requisitos. (¿Qué Es el Backend? Definición y Explicación, 2024)

Para el desarrollo del frontend se aplicarán lenguajes de programación como HTML, CSS y JavaScript, comummente utilizados en desarrollo web, para crear la interfaz de usuario y la experiencia del mismo. En cambio el backend se encargará de la lógica de negocio, la gestión de datos y la comunicación con la base de datos, este se realizará con el lenguaje de codificación de python. (Amazon Web Services, s. f.)

Finalmente la base de datos se diseñará utilizando un modelo relacional, que permite organizar los datos en tablas relacionadas entre sí, facilitando la gestión y consulta de la información. Este modelo es adecuado para el tipo de datos que se manejarán en el sistema, como perfiles de estudiantes, ofertas laborales y postulaciones. Como base de datos especifica se tiene en cuenta PostgreSQL, que es un sistema de gestión de bases de datos relacional de código abierto o MySQL, que es otro sistema de gestión de bases de datos relacional de código abierto, ambos conocidos por su rendimiento y fiabilidad.

Sobre esta infraestructura tecnológica se desarrollará el módulo de recomendación, el cual corresponde a una de las funcionalidades principales del sistema debido a su capacidad para relacionar automáticamente las vacantes disponibles con los perfiles de los estudiantes registrados.

\subsubsection*{Sistema de recomendaciones}
El módulo de recomendación calcula un índice de compatibilidad entre cada oferta publicada y los perfiles de los estudiantes registrados. Este índice se basa en la comparación de atributos clave, como la carrera del estudiante, sus pretensiones salariales, su disponibilidad horaria y su preferencia de modalidad de trabajo (presencial, híbrida o remota). El sistema asigna un peso a cada uno de estos atributos para reflejar su importancia relativa en el proceso de selección. Por ejemplo, la carrera puede tener un peso del 40%, las pretensiones salariales un 30%, la disponibilidad horaria un 20% y la modalidad de trabajo un 10%.
El algoritmo funciona aplicando una fórmula de compatibilidad que combina los pesos asignados a cada atributo con la coincidencia entre los datos de la oferta y el perfil del estudiante. Por ejemplo, si un estudiante tiene una carrera que coincide con la oferta, se suma el peso correspondiente a esa coincidencia. Si las pretensiones salariales del estudiante están dentro del rango ofrecido por la empresa, se suma el peso de esa coincidencia, y así sucesivamente para cada atributo. El resultado final es un índice de compatibilidad que se utiliza para ordenar las recomendaciones en función de su relevancia. Esta lista se actualiza dinámicamente a medida que se publican nuevas ofertas o se actualizan los perfiles, asegurando que las recomendaciones sean siempre relevantes y oportunas.

Para que este módulo funcione de manera integrada con el resto de las funcionalidades de la plataforma, es necesario definir una arquitectura de software que permita organizar adecuadamente los distintos componentes del sistema y asegurar su correcto funcionamiento.

\subsubsection*{Arquitectura y Modelo de calidad}
El sistema se diseñará utilizando una arquitectura de software basada en el patrón de diseño Modelo-Vista-Controlador (MVC), que permite separar la lógica de negocio, la presentación y la gestión de datos en componentes independientes. El modelo se encargará de la gestión de datos y la interacción con la base de datos, la vista se encargará de la presentación de la información al usuario a través del frontend, y el controlador se encargará de manejar las solicitudes del usuario y coordinar la interacción entre el modelo y la vista. Esta arquitectura facilita el mantenimiento y la escalabilidad del sistema, permitiendo que cada componente pueda ser desarrollado y modificado de forma independiente sin afectar a los demás.

En cuanto al modelo de calidad, se aplicará el modelo FURPS+ para la etapa de análisis de requisitos, debido a que ayuda a clasificar y comprender como afectan los requisitos del sistema en la funcionalidad, usabilidad, confiabilidad, rendimiento y soporte del sistema. Además se aplicará el modelo ISO/IEC 25010 para el plan de pruebas, esto porque proporciona un marco para evaluar la calidad del software en términos de seguridad, fiabilidad, capacidad de interacción del usuario con el sistema y las fionalidades, fundamental para garantizar que el sistema cumpla con los estándares de calidad establecidos por la industria y para los usuarios.

Una vez definidos los aspectos relacionados con la arquitectura y la calidad del software, también es importante considerar las estrategias de mantenimiento y soporte necesarias para asegurar la continuidad operativa del sistema y su correcto funcionamiento a largo plazo.

\subsubsection*{Mantenimiento y Soporte}
Es importante tener algunos aspectos en consideración como que el modulo de recomendación es el componente más expuesto a fallas de tipo Crash y Algoritmos, porque cruza múltiples tablas y aplica lógica de compatibilidad, lo que puede generar errores si los datos no están bien estructurados o si hay inconsistencias en la información. Por lo tanto, se consideraría implementar mecanismos de monitoreo y alertas para detectar cualquier anomalía en el funcionamiento del módulo de recomendación, así como pruebas unitarias exhaustivas para garantizar estabilidad y rendimiento. Además, a largo plazo y si el proyecto se lleva a cabo en la realidad como sistema institucional, habría que implementar un plan de mantenimiento que incluya actualizaciones regulares del sistema, revisión de la base de datos y optimización del código para asegurar que el sistema funcione de manera eficiente y sin interrupciones a lo largo del tiempo.

También se pueden generar fallas de recuperación, esto debido al proceso de postulación, el cual involucra múltiples pasos y puede ser propenso a errores si no se manejan adecuadamente las transacciones y la gestión de datos. Para mitigar este riesgo, se podrían implementar mecanismos de recuperación automática en caso de fallas durante el proceso de postulación, como la capacidad de reintentar la operación o restaurar el estado anterior en caso de errores críticos.

Por otra parte, la validación de empresas podría implicar una mitigación de Byzantine, el riesgo se debe a que una empresa malintencionada o con información falsa podría intentar ingresar al sistema, lo que podría comprometer la integridad de la plataforma y la confianza de los usuarios. Para mitigar este riesgo, se podrían implementar medidas de validación rigurosas para las empresas que solicitan ingreso al portal, como la verificación de su identidad y la validación de su información antes de permitir publicar ofertas laborales en la plataforma. 

Finalmente ante cualquier falla y en caso de la perduración del proyecto, se podrían establecer procedimientos para la gestión de incidencias y la resolución de problemas, asegurando que cualquier falla sea abordada de manera rápida y eficiente, minimizando el impacto en los usuarios y la plataforma.



Del Congreso Nacional, B. (s. f.). Biblioteca del Congreso Nacional. www.bcn.cl/leychile. https://www.bcn.cl/leychile/navegar?idNorma=141599
¿Qué es el backend? Definición y explicación. (2024, 7 junio). Ionos Digital Guide. https://www.ionos.com/es-us/digitalguide/paginas-web/creacion-de-paginas-web/que-es-el-backend/ 
Amazon Web Services. (s. f.). ¿Cuál es la diferencia entre el front end y back end en el desarrollo de aplicaciones? Amazon Web Services, Inc. https://aws.amazon.com/es/compare/the-difference-between-frontend-and-backend/


\subsection{Relación con Otros Proyectos}%(cami)
El proyecto "Finis Trabaja" se relaciona con otros proyectos de desarrollo de plataformas de empleo y vinculación laboral, como LinkedIn, Computrabajo, Indeed y la Bolsa Nacional de Empleo (BNE). Estas plataformas ofrecen servicios similares en términos de publicación de ofertas laborales, gestión de perfiles y recomendaciones, pero "Finis Trabaja" se diferencia al estar específicamente diseñada para la comunidad académica de la Universidad Finis Terrae, con un enfoque en conectar a estudiantes y egresados con oportunidades laborales relevantes a sus áreas de estudio. Además, el módulo de recomendación automática basado en atributos específicos del perfil académico es una característica distintiva que busca optimizar la conexión entre los usuarios y las ofertas laborales disponibles.

Adicionalmente, desde una perspectiva transaccional o de intercambio de información, el sistema podría relacionarse con sistemas institucionales y administrativos. Tales como, los sistemas académicos o institucionales de la universidad, estos proporcionan información relevante sobre los estudiantes, como la carrera cursada, estado académico (activo o inactivo), datos de contacto y otros antecedentes académicos. Esta información podría utilizarse para la validación de perfiles, del Currículum Vitae Virtual (CVV) y la gestión de usuarios.

Además, a futuro, podría evolucionar con la integración de mecanismos de autenticación única (SSO) para los usuarios, permitiendo que los estudiantes inicien sesión en la plataforma utilizando sus credenciales institucionales, lo que simplificaría el acceso y mejoraría la experiencia del usuario al no tener que gestionar múltiples cuentas y contraseñas para diferentes sistemas dentro de la universidad.

Por otra parte, el sistema se relacionaría con servicios de correo electrónico institucional. Esta integración tendría como finalidad enviar notificaciones a los usuarios sobre nuevas ofertas laborales, actualizaciones en sus postulaciones o cambios relevantes en su perfil. Sin embargo, la información principal continuaría centralizada dentro de Finis Trabaja, utilizando el correo únicamente como un mecanismo complementario de comunicación.

Otro aspecto a considerar es el almacenamiento de archivos adjuntos, como certificados o documentos relevantes para el perfil del estudiante. En este caso, se podría integrar con servicios de almacenamiento en la nube, para permitir a los usuarios subir y gestionar los documentos de manera segura y accesible desde la plataforma. Esto facilitaría la gestión de archivos pero mantendría la información principal del perfil dentro de Finis Trabaja para el proceso de recomendación y conexión con las ofertas laborales.

Finalmente, en conjunto, estas integraciones y relaciones con otros sistemas y servicios institucionales contribuirían a crear un ecosistema digital más cohesivo y eficiente, mejorando la experiencia del usuario y optimizando la gestión de información relevante para la vinculación laboral de los estudiantes y egresados de la Universidad Finis Terrae.


\section{Requisitos del sistema}%(cami)



\subsection{Lista de Requisitos Reducida}%(cami)



\subsection{Especificación de Requisitos}%(cami)




\subsection{Matriz de Trazado RU/CU ó HU/CU}%(cami)



\end{document}
